\subsection{Variétés allotropiques de l'étain}
% Exercices et problèmes corrigés de thermodynamique chimique, pg2
L’étain (\ce{Sn}) existe sous deux formes allotropiques, l’étain blanc et l’étain gris. Quelle est la forme la plus stable à \SI{25}{\celsius}, sachant que l’entropie molaire standard absolue ($S_m^\circ (298 ~ \si{\kelvin})$) de l’étain blanc est égale à \SI{26.33}{\joule\per\mole\per\kelvin} et que
celle de l’étain gris est égale à \SI{25.75}{\joule\per\mole\per\kelvin} et que la variation de l’enthalpie $\Delta H^\circ (298 ~ \si{kelvin})$ due à la transformation d’étain blanc en étain gris est égale à \SI{2.21}{\kilo\joule\per\mole}.
